\section{Main}\label{sec:introduction}

% \begin{figure}[h!]
%     \centering
%     {{canva:2}}
%     \caption{a. Not all genes that are trait relevant are associated with a genetic variant that rises to significance. b. Perhaps shared genetic annotations can allow us to infer trait relevance from genes that are missed by GWAS through a novel measure of indirect genetic support. c. If indirect genetic support can be put on the same scale as direct genetic support, we can recover relevant genes missed by GWAS.}
%     \label{fig:pigean-intuition}
% \end{figure}

% Broad introductory paragraph
Most human traits are influenced by many genes, and most human genes influence many traits. Mapping this complex web of gene–trait relationships is essential for uncovering disease biology \cite{cohen_sequence_2006}, understanding clinical heterogeneity \cite{udler_type_2018}, accurately diagnosing patients \cite{yang_clinical_2013}, and designing effective therapies \cite{plenge_validating_2013}. Among the most widely used and powerful approaches to to identify gene-trait relationships are genetic association studies. For common traits, genome-wide association studies (GWAS) have identified tens of thousands of trait-associated loci across thousands of phenotypes \cite{abdellaoui_15_2023}, while rare disease studies have linked thousands of Mendelian conditions to specific causal genes \cite{chung_meta-analysis_2023}. A genetic association provides a gene with ``genetic support" for a human trait -- establishing the human-relevance and causal effect of its perturbation.

The number of trait-associated genes identified by genetic association studies is, however, but a small fraction of the total number of trait-relevant genes -- that is, those that can be perturbed in humans to yield an effect on a trait of interest \cite{finan_druggable_2017}. Many trait-relevant genes escape detection (Figure \ref{fig:pigean-intuition-a}) due to limited statistical power \cite{wang_statistical_2019}, constraining evolutionary pressure on genetic variation \cite{gazal_linkage_2017}, or the difficulty of assigning noncoding GWAS signals to their effector genes  \cite{schipper_demystifying_2022,costanzo_realizing_2025}. Yet the full set of trait-relevant genes is of most interest for drug development and experimental study \cite{jhamb_pathway_2019,macnamara_network_2020}. In part due to this incompleteness, most translational research does not use human genetics to prioritize genes \cite{wehling_assessing_2009,dornbos_evaluating_2022}, even with the  increasing recognition that drug targets with genetic support are substantially more likely to succeed clinically \cite{plenge_validating_2013,minikel_refining_2024}. 

% Related Work and More specific stage setting to elicit the the contributions paragraph
By contrast, non-human-genetic approaches are more widely used to identify trait-relevant genes but come with critical tradeoffs. Experimental studies in animal \cite{seok_genomic_2013,widjaja_inhibiting_2019} and cell-based \cite{ben-david_genetic_2018,canon_clinical_2019} systems can establish causal links between genes and phenotypes, but their relevance to humans is often uncertain. Human observational studies, such as gene expression profiling \cite{wu_integrating_2022,rodriguez_machine_2021} or epidemiological associations \cite{barter_cholesteryl_2011,bhatt_cardiovascular_2019}, provide human-relevant context but lack causal directionality. The literature and biomedical databases are rife with gene-trait relationships derived from these or other non-genetic studies, and they are widely used as a starting point for drug discovery or experimental study \cite{swinney_how_2011,dornbos_evaluating_2022} (Figure \ref{fig:pigean-intuition-b}). Recent studies demonstrate how these sources can be rapidly and accurately queried with large language models (LLMs) \cite{toufiq_harnessing_2023,hu_evaluation_2025}. However, it is well documented that prioritizing genes using these evidence sources is prone to bias, incompleteness, or circularity \cite{brito_removing_2011,gillis_assessing_2013,gaudet_removing_2017,dunham_human_2018,haynes_gene_2018,stoeger_large_2018,richardson_meta_2024,tan_caution_2025}.

Ideally, a map of gene-trait relationships would combine the complementary benefits of human genetic and non-human-genetic evidence sources. It is commonplace to evaluate genes individually through manual expert synthesis of genetic data and the published literature \cite{mountjoy_open_2022,karjalainen_genome-wide_2024}, but a quantiative map across the genome and phenome requires scalabale statistical methods. Such methods have been developed both in the context of GWAS effector gene identification, where gene ``functional annotations'' enriched for genetic associations are used to prioritize genes nearby associated SNPs \cite{pers_biological_2015,weeks_leveraging_2023}, and in the context of rare disease diagnostics, where literature mining \cite{birgmeier_amelie_2020} and interaction networks \cite{smedley_next-generation_2015} are used to prioritize genes that carry variants observed in a patient. These approaches, however, rely on heuristics (such as harmonic averaging \cite{buniello_open_2025} or regressions using noisy or incomplete training data \cite{weeks_leveraging_2023,schipper_prioritizing_2025}) and often lack a biological model, which limits their portability across traits or application to settings other than that for which they were originally designed. As a result, the currently charted map of trait-gene relationships across rare and common diseases remains sparse, inaccurate, and inconsistent across traits.

% Contribution explanation
Here, we combine human genetic association data for {{num_common_traits}} common and {{num_rare_traits}} rare traits and non-human-genetic data encoded in {{num_available_gene_annotations}} gene annotations to systemaically identify trait-relevant genes (Figure 1A-C) to form a Genetic Effect Matrix over All Phenotypes (GEMAP). Each gene-trait relationship is a composite of two statistically calibrated values: a traditional measure of  ``direct genetic support'' (the strength of gene-trait association from human genetic data) and a new  measure of ``indirect genetic support'' (the extent to which the gene shares functional annotations with other trait-associated genes). Indirect genetic support is measured on the same scale as direct genetic support, estimates a measure of gene-trait relevance distinct from direct gene-trait association, and enables us to use unbiased human genetic associations to calibrate noisy and heterogeneous non-human-genetic data sources and learn which are trustworthy for each trait. We validated our method on three independent evidence axes—genetic, pharmacologic, and clinical. This map is orders of magnitude denser than manually curated knowledge graphs and can be used to calculate a novel resiliancy over every trait that provides novel insights into the mechansic heterogeneity of traits. 

%This map when thresholded is orders of magnitude denser than manually curated knowledge graphs and, when unthresholded, can be viewed as dual numerical embeddings of genes and traits learned from combined genetic and non-genetic data sources. We illustrate the utility of this resource to aid in the interpretation of gene function and trait mechanistic heterogeneity, predict drug target relative success more accurately and comprehensively than previous resources, and provide a resource for agentic AI not capturable by pretraining on text-only corpora.

%This decomposition formalizes and generalizes previous heuristic or machine learning-based integration strategies \cite{smedley_next-generation_2015,weeks_leveraging_2023,schipper_prioritizing_2025}, offering a statistically grounded alternative.

%Unlike conventional biomedical knowledge graphs \cite{li_graph_2022}, the PGEG avoids lossy edge labeling \cite{himmelstein_systematic_2017,morris_scalable_2023}, heuristic integration \cite{ochoa_next-generation_2023}, and manual curation \cite{putman_monarch_2024}, while remaining directly usable for downstream applications.
